% ============================================================
% FILE RỜI — CHƯA \input VÀO main.tex
% Nội dung: khảo sát nguồn dữ liệu nông sản cho mô-đun AI (FR12.2 dự báo
% nhu cầu/giá, gián tiếp FR12.1 gợi ý sản phẩm).
% Đề xuất vị trí khi merge: chapters/chuong2-kienthuc-nentang.tex, thành một
% \subsection mới trong \section{Cơ sở lý thuyết nền tảng} (sec:lythuyet),
% đặt ngay sau \subsection{Truy xuất nguồn gốc...} (subsec:gs1blockchain)
% và trước \subsection{Hệ thống đánh giá...} (subsec:reputation) — vì đây là
% mạch "dữ liệu nền cho AI" tiếp nối mạch "dữ liệu nền cho truy xuất".
% Sau khi merge: xóa comment này, đổi \subsection*{...} thành \subsection{...}
% kèm \label phù hợp, và bổ sung các @online/@misc dưới đây vào references.bib.
% ============================================================

\subsection*{Nguồn dữ liệu nông sản cho phân tích và mô hình AI}
\label{subsec:nguondulieunongsan}

Hai mô-đun AI trong phạm vi đồ án --- gợi ý sản phẩm (FR12.1) và dự báo nhu cầu/giá theo mùa vụ (FR12.2, thiết kế chi tiết tại mục~\ref{sec:thietkeAI}) --- cần dữ liệu nông sản làm đầu vào huấn luyện và đánh giá. Mục này khảo sát các nguồn dữ liệu công khai tại Việt Nam có thể phục vụ mục đích đó, đối chiếu với dữ liệu nội sinh mà chính nền tảng sẽ tạo ra sau khi vận hành.

% TODO: nhóm xác nhận lại số liệu cụ thể (18.000 vs 18.500 sản phẩm, 24 vs 26 tỉnh ở nguồn thứ 3) trước khi trích dẫn chính thức, vì các bài báo tham khảo có chênh lệch nhỏ theo thời điểm đăng.
Khảo sát nhanh mục ``Nông nghiệp'' trên Cổng dữ liệu mở quốc gia (\url{open.data.gov.vn}) cho thấy chỉ có hai bộ dữ liệu được công bố tại thời điểm khảo sát --- khác biệt rõ rệt so với các lĩnh vực khác trên cùng cổng. Đây là cơ sở cho nhận định rằng nguồn dữ liệu mở chuyên biệt cho nông nghiệp tại Việt Nam còn rất hạn chế, và là một hạn chế cần nêu thẳng khi đánh giá mô-đun AI ở Chương~\ref{chap:danhgia} thay vì né tránh.

Bảng~\ref{tab:nguondulieu} tổng hợp các nguồn đã khảo sát.

\begingroup
\small
\renewcommand{\arraystretch}{1.25}
\begin{longtable}{L{2.6cm} L{3.0cm} L{4.2cm} L{3.2cm}}
\caption{Nguồn dữ liệu nông sản khảo sát cho mô-đun AI}
\label{tab:nguondulieu} \\
\toprule
\rowcolor{tblheadbg}
\textbf{Nguồn} & \textbf{Đơn vị quản lý} & \textbf{Loại dữ liệu} & \textbf{Mức khả dụng cho Farmery} \\
\midrule
\endfirsthead
\toprule
\rowcolor{tblheadbg}
\textbf{Nguồn} & \textbf{Đơn vị quản lý} & \textbf{Loại dữ liệu} & \textbf{Mức khả dụng cho Farmery} \\
\midrule
\endhead
\bottomrule
\endfoot
Cổng dữ liệu mở quốc gia (\texttt{data.gov.vn}/\texttt{open.data.gov.vn}) & Bộ Khoa học và Công nghệ & Đa ngành; mục Nông nghiệp hiện chỉ 2 bộ dữ liệu & Thấp --- chỉ dùng đối chiếu phụ \\
\rowcolor{tblaltbg}
Tổng cục Thống kê (GSO/NSO) & Bộ Tài chính & Diện tích gieo trồng, năng suất, sản lượng theo cây trồng/địa phương/năm & Trung bình --- tần suất năm/quý, phù hợp baseline theo mùa vụ \\
Hệ thống truy xuất nguồn gốc nông sản quốc gia (CheckVN, \texttt{traceviet.mae.gov.vn}, vận hành từ 1/7/2026) & Bộ Nông nghiệp và Môi trường & Batch, nông hộ, vùng trồng, mã số vùng trồng, cơ sở đóng gói & Cao về tham chiếu thiết kế schema; thấp về khả năng lấy dữ liệu huấn luyện trực tiếp \\
\rowcolor{tblaltbg}
Mã số vùng trồng (QĐ 3156/QĐ-BNN-TT, TCCS 774:2020/BVTV) & Cục Trồng trọt và Bảo vệ thực vật & Danh mục vùng trồng đã cấp mã theo tỉnh/loại cây & Trung bình --- đối chiếu/xác thực trường vùng trồng trong ERD \\
Bộ Công Thương / Cục Xuất nhập khẩu & Bộ Công Thương & Kim ngạch xuất khẩu nông sản theo thị trường/mặt hàng theo tháng & Trung bình --- dữ liệu ngữ cảnh thị trường \\
\rowcolor{tblaltbg}
Sàn TMĐT nông sản đối chứng (FoodMap, VIPO Mall, Buudien.vn) & Doanh nghiệp & Giá niêm yết công khai trên sàn & Thấp-trung bình --- không có API chính thức, rủi ro pháp lý nếu crawl quy mô lớn \\
FAOSTAT, World Bank Open Data & FAO, World Bank & Sản lượng, xuất-nhập khẩu theo quốc gia (cấp vĩ mô) & Trung bình --- đối chiếu quốc tế, không đủ chi tiết cho dự báo giá nội địa \\
\rowcolor{tblaltbg}
Dữ liệu nội sinh Farmery & Farmery (tự thu thập) & Lịch sử giá thực tế theo sản phẩm/vùng/thời điểm, clickstream & Cao nhất về độ khớp nhưng chỉ có sau khi vận hành (cold-start) \\
\end{longtable}
\endgroup

\noindent\textbf{Nhận định.} Không có một nguồn mở ``sẵn dùng'' đủ tốt để huấn luyện thẳng mô-đun dự báo giá/nhu cầu ngay từ ngày đầu vận hành. Hướng khả thi kết hợp ba lớp: \textbf{(a)} số liệu vĩ mô GSO/FAOSTAT làm baseline theo mùa vụ, \textbf{(b)} dữ liệu nội sinh của Farmery tích lũy dần sau khi vận hành --- dùng heuristic/seasonal naive ở giai đoạn cold-start rồi chuyển sang mô hình học máy khi đủ dữ liệu, đúng hướng đã chọn cho bài toán NC3 (mục~\ref{subsec:vandenghiencuu}) --- và \textbf{(c)} mã số vùng trồng/hệ thống truy xuất quốc gia dùng để \emph{chuẩn hóa schema định danh}, không phải để lấy dữ liệu huấn luyện, nhất quán với nguyên lý định danh GS1 đã chọn tại mục~\ref{subsec:gs1blockchain}.

% ============================================================
% Nguồn tham khảo cần bổ sung vào references.bib khi merge (ví dụ, cần nhóm
% tự kiểm tra lại URL/ngày truy cập trước khi đưa chính thức):
%
% @online{gso_solieu,
%   author  = {{Tổng cục Thống kê}},
%   title   = {{Số liệu thống kê nông, lâm nghiệp và thủy sản}},
%   year    = {2026},
%   url     = {https://www.gso.gov.vn/so-lieu-thong-ke/},
%   urldate = {2026-09-23}
% }
%
% @online{checkvn_traceviet,
%   author  = {{Bộ Nông nghiệp và Môi trường}},
%   title   = {{Hệ thống truy xuất nguồn gốc nông sản quốc gia (CheckVN)}},
%   year    = {2026},
%   url     = {https://traceviet.mae.gov.vn/},
%   urldate = {2026-09-23}
% }
%
% @misc{qd3156_2022,
%   author = {{Bộ Nông nghiệp và Phát triển nông thôn}},
%   title  = {{Quyết định số 3156/QĐ-BNN-TT về mã số vùng trồng}},
%   year   = {2022}
% }
%
% @online{opendata_govvn,
%   author  = {{Bộ Khoa học và Công nghệ}},
%   title   = {{Cổng dữ liệu mở quốc gia -- nhóm Nông nghiệp}},
%   year    = {2026},
%   url     = {https://data.gov.vn/},
%   urldate = {2026-09-23}
% }
% ============================================================

% ============================================================
% VIỆC CẦN NHÓM XÁC NHẬN TRƯỚC KHI MERGE:
% 1. Vị trí chèn: đề xuất trong chuong2-kienthuc-nentang.tex, ngay sau
%    subsec:gs1blockchain — cần nhóm xác nhận lại hoặc chọn vị trí khác
%    (ví dụ mục~sec:hienthucAI ở Ch.4/Ch.6 để mô tả nguồn dữ liệu THẬT sẽ dùng
%    khi hiện thực, tách khỏi phần khảo sát lý thuyết ở Ch.2).
% 2. Xác nhận số liệu hệ thống CheckVN (18.000 vs 18.500 sản phẩm, 24 vs 26
%    tỉnh) trước khi trích dẫn chính thức trong bảng.
% 3. Thêm các entry references.bib nêu trên (đã kiểm tra chưa trùng khóa với
%    references.bib hiện tại — nhóm tự xác nhận lại URL/ngày truy cập).
% ============================================================
